 \documentclass[12pt ]{article} 
 \usepackage{amssymb}
\usepackage{amsmath}
\usepackage{geometry}
\geometry{
 a4paper,
 total={160mm,240mm}  
  }
  
 \begin{document}
 $~$ 
 
 \vspace{2cm}
 
 \centerline{\bf \huge   Naslov}
 
 \vspace{1cm}
 
  \centerline{\huge  Avtor}
  
   \vspace{1cm}
  
  \centerline{\large Mentor}
  
  \vspace{1cm}
  
  \centerline{\large Seminar, 3. letnik}
  
  \vspace{1cm}
  
  \centerline{\large Oddelek za fiziko, FMF, UL, 2021/2022}
  
  \vspace{6cm}
  
  \begin{minipage}[c]{0.9\hsize}
  {\bf Povzetek}
  
  Pri pripravi seminarskih nalog se nau\v cimo zbiranja in kriti\v cne presoje strokovne literature, pisanja strokovnih sestavkov in predstavljanja prav tak\v snih tem v javnosti. Pisanje besedila je zelo pomemben del tega procesa, saj s tem poka\v zemo, ali temo zares  obvladamo.    
    \end{minipage}

 
 \newpage
  
  
 \section{Uvod}
 
 \section{Jezik}
 
       Potrudimo se, da bomo seminarske naloge napisali v lepi in gladko berljivi ter slovni\v cno neopore\v cni zborni sloven\v s\v cini. Pazimo na pravopis in tipkarske napake. Kjer je mogo\v ce, tujko nadomestimo s slovensko ustreznico, obenem pa skrbimo, da uporabljamo ustaljene strokovne izraze.    
    \end{document} 